\documentclass[11pt]{article}

\usepackage[margin=1in]{geometry}
\usepackage[T1]{fontenc}
\usepackage[utf8]{inputenc}
\usepackage{lmodern}
\usepackage{setspace}
\usepackage{hyperref}
\usepackage{booktabs}
\usepackage{array}
\usepackage{longtable}
\usepackage{enumitem}
\usepackage{graphicx}
\usepackage{float}

\hypersetup{
    colorlinks=true,
    linkcolor=black,
    citecolor=black,
    urlcolor=blue,
    pdftitle={ChainSentry Project Report}
}

\setstretch{1.15}

\title{ChainSentry: An Explainable Pre-Signature Transaction Risk Analysis Prototype for Blockchain Users}
\author{Yiheng Lu (123090391) \and Zhuokai Chen (123090057)}
\date{April 2026}

\begin{document}

\maketitle

\begin{abstract}
Blockchain wallet users are routinely asked to sign transactions whose practical consequences are difficult to infer from raw calldata, token amounts, and destination addresses alone. This usability gap becomes a security problem when users approve high-risk actions such as unlimited token allowances, suspicious contract interactions, or permission-changing calls without a clear explanation of what the transaction will actually enable. ChainSentry is a student-built prototype that addresses this gap through a pre-signature risk analysis workflow designed around explainability rather than pure detection performance. The system combines a React-based wallet-facing interface, a FastAPI analysis service, heuristic transaction parsing, lightweight simulation, a local flagged-address detector, and a packaged relation-aware graph model.

The core design choice in ChainSentry is to analyze each pending transaction as a transaction-centered local heterogeneous graph instead of as a full blockchain-scale evolving graph. This design follows directly from the product goal: the system must react to one imminent transaction and explain the result in user-readable language before signing. The graph model therefore operates alongside deterministic detectors and simulation summaries rather than replacing them. The backend returns a concise risk report structured around expected action, risk reason, possible impact, recommendation, and severity. The same structure is then rendered in a front-end risk card format that supports both demonstration scenarios and manual transaction entry.

Evaluation in the repository now centers on automated API tests and graph-pipeline unit tests that exercise the shipped analysis path. These checks show that the system distinguishes clean transfers from high-risk approvals, flagged destinations, operator-control cases, and privilege-escalation scenarios while preserving a clear low-risk control case. Overall, ChainSentry shows that a narrow, explainable, transaction-centered hybrid architecture is a practical direction for student-scale blockchain security tooling.
\end{abstract}

\tableofcontents
\newpage

\section{Introduction}

Wallet confirmation interfaces typically expose low-level transaction details such as destination address, token amount, calldata, gas cost, and method name. For expert users, this information may be enough to infer intent. For non-expert and intermediate users, however, the same interface is often too technical to interpret quickly and reliably. A transaction that appears to be a harmless swap may actually grant a persistent allowance. A call that looks like routine asset management may in practice transfer collection-wide operator control. A transfer may route value to an already flagged destination. The problem is therefore not only malicious behavior on-chain, but also the mismatch between what a user sees before signing and what the signed transaction can later enable.

ChainSentry was built around this practical problem. The project goal is not to create a full blockchain threat intelligence platform, nor to solve arbitrary smart contract security analysis at web scale. Instead, the project focuses on a narrow but important setting: pre-signature transaction assistance for users who need fast, understandable guidance. The prototype accepts a pending transaction, normalizes its fields, performs targeted risk analysis, and returns a short report that states what the transaction is expected to do, why it may be risky, what could happen after signing, and what action the user should take next.

This design objective has several implications. First, the system must remain lightweight enough to run in a demo or local deployment environment. Second, the output must be explanation-oriented rather than a bare score. Third, the model architecture should match the product surface. Since the prototype analyzes one pending transaction at a time, a local transaction-centered graph is more appropriate than a global streaming graph over the full chain. Fourth, the machine learning component should complement deterministic checks instead of replacing them. In practice, some risk patterns, such as a flagged destination or an unusually large approval, are easier to explain when a direct detector exists.

Based on these principles, ChainSentry adopts a hybrid architecture. The front end is built with React and Vite and provides both wallet-assisted input and pre-filled demonstration scenarios. The backend is built with FastAPI and exposes a stable analysis endpoint. Internally, the backend parses incoming transaction data, derives a normalized transaction type, runs a heuristic simulation engine, constructs a local heterogeneous transaction graph, extracts scalar features, and applies a relation-aware graph model with multiple output heads. Heuristic findings and graph-model signals are then fused into a unified risk report. The response is intentionally structured to support a user-facing risk card instead of a security analyst dashboard.

The remainder of this report is organized as follows. Section 2 reviews related work on Ethereum phishing detection, graph-based blockchain anomaly detection, and transaction simulation. Section 3 describes the system architecture. Section 4 explains the workflow and design decisions. Section 5 gives implementation details grounded in the repository. Section 6 summarizes evaluation results from the current automated tests and end-to-end scenarios. Section 7 discusses limitations and future work. Section 8 provides the required equal-contribution statement.

\section{Related Work and Background}

\subsection{Blockchain Risk Detection as a Graph Problem}

Blockchain activity is naturally relational. Addresses interact with contracts, contracts move tokens, transactions propagate permissions, and suspicious behavior may only become visible when neighboring entities and interaction patterns are considered together. For that reason, graph-based methods have become increasingly common in blockchain security research. Recent phishing- and anomaly-detection studies repeatedly treat transaction data as a graph problem because flat tabular approaches struggle to capture the dependence between users, contracts, and fund flows \cite{chen2024dahgnn,chang2025anomalous}.

The same intuition appears in anomaly-detection work outside the exact phishing setting. Chang et al. propose SGAT-BC, a graph-neural anomaly detector that combines subtree attention, graph attention, and ensemble learning for imbalanced blockchain datasets \cite{chang2025anomalous}. Their work is especially relevant because it foregrounds two issues that also matter in ChainSentry: graph structure is useful for fraud analysis, and class imbalance remains a practical obstacle. Even when graph models achieve strong metrics, the quality of labels, the rarity of malicious events, and the dependence on dataset construction strongly shape the final system behavior.

ChainSentry adopts the graph perspective but differs in scope. Most published anomaly-detection systems aim at address-level or node-level classification over larger transaction graphs. By contrast, ChainSentry centers the graph on a single pending transaction so that its output can be consumed before signature. This local-graph choice sacrifices large-scale relational coverage, but it improves alignment with the actual user task: should this specific transaction be signed now?

\subsection{Ethereum Phishing and Scam Detection}

Ethereum phishing detection has evolved from manual feature engineering toward graph-oriented and hybrid deep learning methods. Earlier systems often relied on transaction statistics such as degree counts, amounts, and temporal summaries. Later work learned embeddings from transaction networks or subgraphs. Chen et al. survey several of these directions and categorize the literature into manual feature extraction, embedding-based methods, and hybrid feature-fusion models \cite{chen2024dahgnn}. Their DA-HGNN system combines data augmentation, temporal modeling, and graph reconstruction, and reports strong phishing-detection performance on curated datasets.

The project is also closely related to more recent phishing research that explicitly studies the role of feature choice. The arXiv paper \emph{Fishing for Phishers: Learning-Based Phishing Detection in Ethereum Transactions} examines the effectiveness of explicit transactional features versus implicit graph-based features and highlights the importance of feature selection, class imbalance, and dataset composition in phishing detection \cite{alghuried2025fishing}. This is directly relevant to ChainSentry because the current prototype intentionally combines graph structure with scalar transaction features instead of depending on a purely graph-only representation.

Industry-facing research makes a similar point from a deployment angle. The GoPlus Security open research note on phishing address detection describes a spatio-temporal fusion network that combines transaction subgraphs with ordered transaction sequences and reports improvements over handcrafted and graph-only baselines \cite{goplus2025phishing}. Although the GoPlus setting is address classification rather than pre-signature assistance, it reinforces an assumption that also holds for this project: phishing and scam behavior are difficult to capture from a single flat feature view, and both relational and temporal context matter.

ChainSentry differs from these systems in two ways. First, it operates on transaction requests rather than retrospective address histories alone. Second, it is designed to emit human-readable findings. Many research systems optimize phishing detection performance but do not specify how their outputs become short user explanations at the signing surface. ChainSentry treats that translation layer as part of the system, not as an afterthought.

\subsection{Smart Contract Fraud and Heterogeneous Graph Modeling}

Related work also exists at the smart contract and heterogeneous-network level. Kanezashi et al. explore Ethereum fraud detection with heterogeneous graph neural networks and argue that typed entities and typed relations can improve fraud modeling over simpler graph abstractions \cite{kanezashi2022hetero}. The core lesson for ChainSentry is not that a student prototype needs a large heterogeneous model, but that homogeneous graph abstractions can miss distinctions that matter when addresses, contracts, tokens, and effects play different roles.

This observation supports one of ChainSentry's main design decisions: the graph is heterogeneous even though it is local. The backend distinguishes transaction, address, contract, token, and effect nodes, and it preserves relation types such as \texttt{initiates}, \texttt{targets}, \texttt{approves}, \texttt{requests\_allowance\_for}, and \texttt{triggers\_effect}. This is a smaller and simpler graph than those used in semantic contract-classification research, but it encodes the same principle that typed relations matter.

\subsection{Transaction Simulation and Wallet Security}

Graph-based scoring alone does not fully solve the user problem at signing time. Some risks are best explained by anticipated post-signature effects. Transaction simulation is therefore an important complementary idea in wallet security. The Zengo transaction simulation white paper, produced under an Ethereum Foundation grant, argues that transaction simulation is becoming a necessary security capability for modern wallets, but also warns that simulation is not a complete defense and must be combined with additional mechanisms such as reputation signals and context-aware explanation \cite{beery2025simulation}. Its recommendations are especially relevant for approval flows: simulation results should explain the implications of approvals and include information about spender reputation.

ChainSentry follows this layered view. The simulation component is intentionally lightweight and heuristic in the current prototype, yet it exists to expose effects that are easier to communicate as consequences than as raw model probabilities. For example, the backend can express that an operator receives broad transfer authority or that a downstream outflow becomes possible after approval. This choice echoes the white paper's broader argument that user protection requires context, not just raw simulation traces.

\subsection{Background Summary and Project Positioning}

The literature suggests three recurring themes. First, graph structure is useful for blockchain security because malicious behavior is relational. Second, phishing and fraud detection often improve when structural and temporal information are fused. Third, end-user protection at signing time also requires consequence modeling and explanation. ChainSentry positions itself at the intersection of these themes. It is not a large-scale detector for all blockchain fraud, but an explainable pre-signature assistant that uses graph reasoning, deterministic checks, and simulation-informed explanation in a single workflow.

\section{System Architecture}

ChainSentry implements a hybrid client-server architecture with five main layers: a front-end interface, an API layer, a transaction normalization and analysis layer, a graph-model inference layer, and an explanation layer. The components and responsibilities are summarized in Table~\ref{tab:architecture}.

\begin{table}[H]
\centering
\caption{Major ChainSentry components and responsibilities.}
\label{tab:architecture}
\begin{tabular}{p{3cm} p{10.5cm}}
\toprule
Component & Responsibility \\
\midrule
React front end & Captures transaction fields, supports wallet-assisted address filling, provides pre-built demo scenarios, submits analysis requests, and renders risk reports as user-facing cards. \\
FastAPI API & Exposes \texttt{POST /api/v1/analyze} and returns a typed analysis response containing normalized transaction data, severity, recommended action, findings, and simulation summary. \\
Parser and detectors & Normalize raw transaction input into a typed transaction kind, infer method selectors when available, detect large or unlimited approvals, check locally flagged destinations, and interpret simulated downstream effects. \\
Simulation and graph ML & Run heuristic simulation effects, build a local heterogeneous transaction graph, extract scalar features, and apply a relation-aware graph model with multiple risk heads. \\
Explanation layer & Merge heuristic findings and model evidence into concise explanations using a fixed structure: expected action, risk reason, possible impact, recommended action, and severity. \\
\bottomrule
\end{tabular}
\end{table}

At the interface level, the front end is implemented with React 19, Vite, Wagmi, and Viem. It exposes fields for chain ID, sender and destination addresses, method name, token amount, approval amount, spender, contract label, calldata, notes, and simulation profile. It also includes five sample scenarios that correspond to meaningful security cases in the backend tests: a clean transfer baseline, a large approval to a benign router, a transfer to a flagged destination, an operator-control approval, and a privilege-escalation call.

At the service level, FastAPI receives a typed \texttt{TransactionRequest}, validates formats such as Ethereum addresses and hex calldata, and routes the request to a central \texttt{analyze\_transaction} function. This function is the main orchestration point of the backend. It invokes the parser, simulation engine, graph-model predictor, severity summarizer, and recommendation mapper before returning the final response.

The backend analysis surface is intentionally narrow. The parser infers one of several transaction kinds: approval, transfer, swap, privilege, native transfer, or generic contract call. The detector layer checks for large or unlimited approvals and for addresses appearing in a local flagged-contract dataset. The simulation engine adds effect summaries such as allowance grant, operator control, privilege escalation, or unexpected outflow. The graph-model inference layer uses this normalized and simulated transaction state to construct a local graph and to compute category-specific risk scores. The explanation layer finally combines these ingredients into user-readable findings.

This architecture supports two important design goals. First, it allows each risk finding to be tied to an identifiable source of evidence, whether a deterministic rule, a simulation effect, or model-generated graph evidence. Second, it allows the system to remain robust if the graph model is unavailable. The inference module includes a heuristic fallback path, so the system can still produce explainable findings even when the learned artifact is missing or cannot be loaded.

\section{Workflow and Design Description}

\subsection{End-to-End Workflow}

The ChainSentry workflow can be described as a seven-stage pipeline:

\begin{enumerate}[leftmargin=2em]
    \item The user enters or loads a pending transaction in the front end.
    \item The front end validates obvious format errors and submits the request to the backend analysis endpoint.
    \item The backend normalizes the raw request into a typed transaction object with an inferred method, transaction kind, selector, and natural-language summary.
    \item A simulation stage derives high-level effects such as reusable allowance, operator-wide control, privilege grants, or unexpected downstream outflow.
    \item The graph layer builds a transaction-centered heterogeneous graph and extracts auxiliary scalar features.
    \item The predictor combines graph-model outputs with heuristic findings and transforms them into final risk findings.
    \item The front end renders a risk report containing overall severity, recommended action, findings, and supporting evidence.
\end{enumerate}

Although this workflow is simple, its structure reflects the main thesis of the project: the system should help a user understand one imminent transaction before they sign it. This is why the pipeline is transaction-centered instead of chain-centered, and why natural-language summarization is built into the response schema rather than deferred to the client.

\subsection{Transaction-Centered Graph Design}

The graph representation is one of the most important design decisions in the project. Instead of constructing a broad historical graph for every address, the backend creates a compact local graph around the current transaction. This graph is anchored by a transaction node and then expanded with typed neighboring entities. The checked-in implementation uses five node types and ten edge types, as shown in Table~\ref{tab:graphschema}.

\begin{table}[H]
\centering
\caption{Node and edge types used in the transaction-centered graph.}
\label{tab:graphschema}
\begin{tabular}{p{4cm} p{8.8cm}}
\toprule
Element type & Members in the current implementation \\
\midrule
Node types & transaction, address, contract, token, effect \\
Edge types & initiates, targets, approves, requests\_allowance\_for, transfers\_value\_to, transfers\_token\_to, routes\_to, grants\_operator\_to, grants\_privilege\_to, triggers\_effect \\
\bottomrule
\end{tabular}
\end{table}

This design preserves the typed structure of blockchain interactions without requiring a large graph database or expensive neighborhood expansion during inference. For example, an approval transaction creates a transaction node, an initiator address node, a destination or spender contract node, a token node, and any simulation effect nodes generated by the simulation engine. The graph thus captures not only who interacts with whom, but also what kind of relation connects them.

The local graph also fits the explainability goal better than a fully global graph. Because the model sees only the entities and effects surrounding one transaction, the resulting evidence can be summarized in a compact way, for example by reporting the graph size and the category-specific score associated with approval, destination, or simulation risk.

\subsection{Hybrid Analysis Instead of Pure ML}

Another key design choice is the use of hybrid analysis. ChainSentry does not treat the graph model as the sole source of truth. Instead, heuristic detectors and heuristic simulation are treated as first-class analysis sources. This is a pragmatic choice for a student system. Certain conditions, such as whether an approval amount exceeds a large or effectively unlimited threshold, are directly observable and easier to explain deterministically. Similarly, a local flagged-address list gives a precise reason for a destination warning. The graph model then acts as an additional relational signal rather than a black-box replacement.

This hybrid design has three benefits. First, it improves reliability in small-data settings. Second, it supports graceful degradation when the model is unavailable. Third, it produces more interpretable outputs because a finding can carry both rule-based evidence and model evidence. In the checked-in inference logic, if a heuristic finding already exists for a risk category, the backend appends graph-model evidence to that finding instead of replacing it. If no heuristic finding exists but the model score crosses a threshold, the backend can emit a generic model-based finding for the relevant category.

\subsection{Explanation-Centered Output Design}

The response schema is designed around user interpretation rather than raw probability reporting. Each finding contains an identifier, category, severity, expected action, risk reason, possible impact, recommended action, and evidence list. This structure appears consistently across approval, destination, and simulation findings. The front end then renders these fields in a card layout, together with the overall severity and a recommendation such as \texttt{proceed}, \texttt{inspect\_further}, or \texttt{reject}.

The importance of this design should not be underestimated. Many blockchain security tools can detect suspicious patterns, but ordinary users still face a translation problem: what does the warning mean in plain language, and what should they do now? ChainSentry treats that translation as a primary system requirement. In that sense, the explanation layer is not just presentation logic; it is part of the product's security value.

\section{Implementation Details}

\subsection{Typed API and Transaction Normalization}

The backend models are implemented with Pydantic. The incoming \texttt{TransactionRequest} includes fields such as \texttt{chain\_id}, \texttt{from\_address}, \texttt{to\_address}, \texttt{method\_name}, \texttt{calldata}, \texttt{value\_eth}, \texttt{token\_symbol}, \texttt{token\_amount}, \texttt{approval\_amount}, \texttt{spender\_address}, \texttt{contract\_name}, and \texttt{simulation\_profile}. Validators normalize addresses to lowercase, enforce hexadecimal address and calldata formats, and normalize optional textual fields.

The parser converts this request into a \texttt{NormalizedTransaction}. It first attempts to derive a method from the calldata selector and then falls back to explicit method name or heuristic inference. A transaction kind is derived from the normalized method and transaction value. The parser then produces a concise natural-language summary such as ``Approve 0x1234...abcd to spend unlimited USDC'' or ``Grant elevated permissions through Shadow Vault.'' This summary is reused throughout the explanation layer, which helps keep the reporting consistent.

\subsection{Heuristic Detectors and Simulation Layer}

The detector layer implements three direct checks: approval-risk detection, flagged-destination detection, and simulation-effect interpretation. Approval detection uses configured thresholds for large and effectively unlimited approvals. Flagged-destination detection reads from a local JSON dataset of demo-flagged addresses. Simulation-based findings are derived from the output of the heuristic simulation engine.

The simulation engine is intentionally simple. It does not execute a full on-chain trace. Instead, it maps transaction kinds and simulation profiles to effect summaries that represent likely consequences relevant to user understanding. For example, an approval with a nonzero approval amount yields an \texttt{allowance\_grant} effect stating that the spender receives reusable token allowance. \texttt{setApprovalForAll} yields an \texttt{operator\_control} effect. A privilege-escalation profile yields a \texttt{privilege\_grant} effect. An \texttt{unexpected\_outflow} profile indicates broader post-signature asset movement.

This layer is modest compared with industrial simulation infrastructure, but it is effective for the prototype's goal. It surfaces exactly the kinds of outcomes that a user often fails to infer from wallet UI alone.

\subsection{Graph Construction and Feature Extraction}

After parsing and simulation, the backend constructs a \texttt{TransactionGraph}. Each node carries a node type and optional attributes. The transaction node includes attributes such as chain ID, transaction kind, method name, and selector. Address, contract, token, and effect nodes carry role-appropriate metadata. Edges encode the typed relations among these nodes.

In parallel, the backend extracts a \texttt{ScalarFeatureSet} with numeric, categorical, and boolean fields. The current numeric features are \texttt{chain\_id}, \texttt{value\_eth}, \texttt{token\_amount}, \texttt{approval\_amount}, \texttt{graph\_node\_count}, \texttt{graph\_edge\_count}, and \texttt{simulation\_effect\_count}. The categorical features are \texttt{transaction\_kind}, \texttt{method\_name}, \texttt{selector}, and \texttt{simulation\_profile}. The boolean features indicate whether a spender exists, whether a token exists, whether simulation was triggered, and whether the request contains a selector.

An important implementation detail is that the current artifact focuses on structural typing and scalar context rather than direct address identity. The repository documentation states that the current vectorization path does not directly encode raw address strings, contract names, or free-text notes into learned node representations. This keeps the model lightweight and easier to generalize across synthetic and adaptor-built examples, but it also limits how much entity-specific nuance the model can capture.

\subsection{Relation-Aware Multi-Head Graph Model}

The core learned component is the \texttt{RelationAwareGraphModel}. The implementation uses trainable node-type embeddings, per-layer self transformations, and relation-specific linear transformations over edge types. During forward propagation, messages are aggregated relation by relation using the graph's typed edges. The model then combines four information sources into a joint representation: the anchor transaction-node state, the graph-wide mean-pooled state, the graph-wide max-pooled state, and an encoded vector of scalar features.

The default model configuration stored in the repository uses hidden dimension 64, two relation layers, categorical embedding dimension 8, feature hidden dimension 48, head hidden dimension 64, and dropout 0.15. The model has three main binary heads for \texttt{approval}, \texttt{destination}, and \texttt{simulation} risk, plus a multiclass severity head with labels \texttt{low}, \texttt{medium}, \texttt{high}, and \texttt{critical}. The artifact schema also supports auxiliary heads such as \texttt{address\_malicious} and \texttt{failure\_aux}, although those are not directly exposed in the user-facing API.

This architecture is a good fit for the project. It is expressive enough to encode typed relations among transaction actors and effects, while still simple enough to train and ship in a student repository. The model is not a full heterogeneous graph transformer over chain-scale data, but it does implement relation-aware message passing that suits a compact local graph.

\subsection{Inference-Time Fusion and Severity Mapping}

At inference time, the predictor first attempts to load the packaged model artifact. If the artifact is unavailable or fails to load, the backend falls back to heuristic mode instead of regenerating a model locally. When the graph artifact is present, the predictor computes category probabilities and a severity distribution. The backend then merges model scores with heuristic findings by category. If a heuristic finding is already present, the model's graph evidence is appended to that finding. If a heuristic finding is absent but a category score exceeds its threshold, the system creates a generic model-based finding for that category, provided the inferred severity is above \texttt{low}.

The final overall severity is simply the maximum finding severity. The overall recommendation then follows a policy mapping: \texttt{low} maps to \texttt{proceed}, \texttt{medium} maps to \texttt{inspect\_further}, and \texttt{high} or \texttt{critical} map to \texttt{reject}. This mapping is intentionally conservative and easy for users to interpret.

\subsection{Training Data and Packaged Handoff Scope}

During development, the model-building work drew on multiple external datasets rather than only on the compact files bundled with the final handoff. The full data inventory consisted of 15,259 labeled malicious Forta records, 115,436 labeled ETH-Labels addresses, 6,912 EtherScamDB scam records, 9,576 PTXPhish phishing samples plus 300 initial-address records, and 99,999 RAVEN finetuning samples together with 20,000 evaluation samples. These full datasets were not all committed into the public repository because of packaging and ignore-file constraints.

For the public handoff, the repository keeps the runtime inference stack and the packaged graph-model artifact while omitting most of the training workspace and raw external corpora. This keeps the submission smaller and easier to understand, but it also means the repository should be read primarily as an inference-oriented prototype rather than as a full reproducible training release.

\subsection{Front-End User Experience}

The front end reflects the project's explanation-first philosophy. Users can load predefined demo scenarios or manually edit transaction fields. Wallet connectivity is supported through Wagmi, and the connected address can auto-fill the sender field. The client validates addresses, calldata, and nonnegative numeric values before submission. The API client uses a 15-second timeout and returns readable error messages when the backend is unavailable or slow.

Once the analysis returns, the risk report panel shows the overall severity, the recommended action, the normalized expected action, detailed findings, simulation details, and transaction metadata. If no finding exists, the UI explicitly states that no configured approval, destination, or simulation-based risk rules were triggered. This behavior is important for a demo setting because it preserves a clear low-risk control case rather than always producing warnings.

\section{Evaluation and Results}

\subsection{Evaluation Setup in the Repository}

The repository now provides two complementary evaluation surfaces. First, end-to-end API tests validate that the full backend returns the expected findings and severity for representative transaction scenarios. Second, unit tests verify graph construction, scalar feature extraction, and predictor selection in the shipped runtime path.

This evaluation strategy is appropriate for a prototype. It does not constitute a large controlled user study or a production-grade benchmark campaign, but it does verify that the implemented pipeline functions coherently from input schema to user-facing output.

\subsection{Scenario-Based End-to-End Results}

The API tests are particularly valuable because they evaluate the system the way a user would experience it: by submitting a transaction and inspecting the final risk report. Table~\ref{tab:scenarioresults} summarizes the key checked scenarios.

\begin{table}[H]
\centering
\caption{Representative end-to-end scenario outcomes from backend tests.}
\label{tab:scenarioresults}
\begin{tabular}{p{4cm} p{2cm} p{2.4cm} p{6cm}}
\toprule
Scenario & Severity & Recommendation & Main observed findings \\
\midrule
Standard transfer to known wallet & low & proceed & No findings; serves as a low-risk control case. \\
Large approval to benign router & medium & inspect\_further & Large approval plus reusable allowance finding. \\
Unlimited approval to flagged contract & critical & reject & Approval, destination, and simulation categories all triggered. \\
Operator-control approval & critical & reject & Broad operator control and unexpected outflow effects. \\
Privilege escalation via \texttt{grantRole} & high & reject & Privilege-escalation simulation finding. \\
\bottomrule
\end{tabular}
\end{table}

These results align well with the intended product behavior. The system does not over-warn on the clean control case. It escalates a large but non-blacklisted approval to medium severity rather than immediately treating it as catastrophic. It marks an unlimited approval to a flagged destination as critical and collects evidence from multiple risk categories. It also handles non-approval permission changes, which broadens the scope beyond the common approval-drain example.

The remaining unit tests further confirm that the graph builder creates transaction-centered graphs with the expected node and edge types, that scalar features include graph context such as node and edge counts, and that the predictor selection logic safely preserves a heuristic fallback when the packaged graph model cannot be used.

\subsection{Interpretation of Results}

Taken together, the evaluation supports three main conclusions.

First, the prototype achieves functional correctness on its intended scope. The full stack from UI request shape to backend response behaves coherently on the key demo scenarios.

Second, the hybrid design is justified. The system is strongest when deterministic rules, simulation effects, and graph scores support each other. For example, the most convincing critical case in the tests is not driven by a bare model score alone; it is driven by the convergence of approval, flagged destination, and simulation evidence.

Third, the graph-model component is promising but still preliminary. The shipped artifact demonstrates that relation-aware graph modeling over local transaction graphs is viable in this setting, but the public handoff is now inference-oriented rather than a full reproducible training release. That means the repository supports qualitative and end-to-end validation more strongly than broad quantitative claims about model generalization.

\section{Limitations and Future Work}

ChainSentry is a focused prototype, and its limitations are significant but clear.

The first limitation is data scale and label quality. During development, the model work drew on several external datasets, but the public repository does not ship a full reproducible training release or the raw corpora needed to re-run the entire training workflow from scratch. This makes it harder to audit label balance, severity coverage, and quantitative generalization directly from the handoff package. Future work should prioritize a larger curated dataset, more consistent severity annotation, stronger validation splits, and a publishable training bundle.

The second limitation is the simplicity of the simulation layer. The current simulation engine is heuristic rather than a full execution environment. It captures useful user-facing consequences, but it does not model arbitrary EVM state transitions, slippage behavior, or protocol-specific side effects. A more advanced version could integrate a real transaction simulator or a forked-chain execution backend for selected supported transaction families.

The third limitation is representational. The current vectorization path emphasizes node and edge types plus scalar features, but does not directly encode raw address identity, contract names, interaction labels, or free-text notes into the model's learned graph state. This makes the model more generic and easier to train in a student environment, but it reduces its ability to exploit richer entity semantics. Future work could add address embeddings, contract clustering, token reputation features, and protocol labels, while still preserving explainability.

The fourth limitation is evaluation scope. The repository has useful automated tests, but it does not include a formal user study or latency benchmark for interactive deployment. Since the system's purpose is user assistance, future work should measure whether the generated explanations actually improve user decisions, confidence, or error detection compared with raw wallet prompts or score-only warnings.

The fifth limitation is coverage. ChainSentry currently targets a small set of explainable risk patterns: large approvals, unlimited approvals, flagged destinations, reusable allowances, operator control, privilege escalation, and unexpected outflow. It does not attempt full smart contract auditing, generalized scam family discovery, or multi-chain production deployment. This narrowness is deliberate, but it also defines the ceiling of the current prototype.

Several future directions follow naturally from these limitations:

\begin{enumerate}[leftmargin=2em]
    \item Expand the labeled dataset and retrain with the \texttt{standard} or \texttt{full} profile rather than the current quick artifact.
    \item Replace heuristic simulation for selected flows with true EVM transaction simulation and state-diff extraction.
    \item Incorporate richer reputation sources such as protocol labels, address poisoning indicators, and phishing-address intelligence feeds.
    \item Add explicit explanation attribution so that each finding can point to the most influential graph relations or scalar features.
    \item Conduct user evaluation on explanation length, wording, trust calibration, and decision quality.
    \item Benchmark inference latency end to end, including graph building, feature encoding, model scoring, and front-end rendering.
\end{enumerate}

Even with these limitations, the prototype already demonstrates a useful systems insight: in the pre-signature setting, narrow explainable coverage may be more valuable than broad opaque detection.

\section{Member Contribution Statement}

This project was completed with a 50/50 contribution split.

Yiheng Lu (123090391) contributed 50\% of the project effort and was primarily responsible for the Student 2 workstream. This included the graph-model pipeline, dataset adaptation and training workflow, model integration into the backend inference path, and the associated evaluation and documentation work.

Zhuokai Chen (123090057) contributed 50\% of the project effort and was primarily responsible for the Student 1 workstream. This included the front-end and back-end prototype foundation, transaction parsing and heuristic analysis flow, simulation and explanation pipeline, interface integration, and the associated testing and documentation work.

The final prototype, repository state, and written report were jointly reviewed and approved by both members. The team therefore declares equal contribution overall.

\section{Conclusion}

ChainSentry addresses a concrete blockchain usability-security problem: users often sign transactions without understanding their real consequences. The project responds with a transaction-centered, explainable, hybrid architecture that combines deterministic checks, lightweight simulation, and relation-aware graph learning. The result is not a full blockchain security platform, but a practical prototype that produces concise, structured, and user-readable risk reports before signature.

From an engineering perspective, the most important contribution of the project is the close alignment between the model design and the product requirement. Because the system's task is to help users decide on one imminent transaction, a local heterogeneous graph is more appropriate than a chain-scale global graph. Because users need explanations rather than scores alone, heuristic findings and simulation summaries remain central even when a learned model is present. Because the project is student-scoped, graceful fallback behavior and compact training artifacts are treated as practical design constraints.

The evaluation available in the repository is necessarily limited, but it supports the overall thesis. ChainSentry distinguishes low-risk from high-risk demo cases, surfaces overlapping evidence in serious scenarios, and demonstrates a functioning graph-model inference pipeline integrated with deterministic checks and simulation cues. Future work should strengthen the data, deepen the simulation, and add user-centered evaluation. Even in its current form, however, ChainSentry shows that explainable pre-signature assistance is a viable direction for blockchain security tooling. The full codebase for the project is publicly accessible at \url{https://github.com/CodyOutcast/ChainSentry}.

\begin{thebibliography}{99}

\bibitem{chainalysis2024crime}
Chainalysis Team, ``2024 Crypto Crime Trends: Illicit Activity Down as Scamming and Stolen Funds Fall, But Ransomware and Darknet Markets See Growth,'' Jan. 18, 2024. Available: \url{https://www.chainalysis.com/blog/2024-crypto-crime-report-introduction/}. Accessed: April 20, 2026.

\bibitem{chen2024dahgnn}
Z. Chen, S.-Z. Liu, J. Huang, Y.-H. Xiu, H. Zhang, and H.-X. Long, ``Ethereum Phishing Scam Detection Based on Data Augmentation Method and Hybrid Graph Neural Network Model,'' \emph{Sensors}, vol. 24, no. 12, p. 4022, 2024. doi: 10.3390/s24124022.

\bibitem{chang2025anomalous}
Z. Chang, Y. Cai, X. F. Liu, Z. Xie, Y. Liu, and Q. Zhan, ``Anomalous Node Detection in Blockchain Networks Based on Graph Neural Networks,'' \emph{Sensors}, vol. 25, no. 1, p. 1, 2025. doi: 10.3390/s25010001.

\bibitem{alghuried2025fishing}
A. Alghuried, A. Alghamdi, A. Alkinoon, S. Choi, M. Mohaisen, and D. Mohaisen, ``Fishing for Phishers: Learning-Based Phishing Detection in Ethereum Transactions,'' arXiv:2504.17953, 2025. Available: \url{https://arxiv.org/abs/2504.17953}. Accessed: April 20, 2026.

\bibitem{goplus2025phishing}
GoPlus Security, ``Phishing Address Detection,'' GoPlus Open Research, 2025. Available: \url{https://whitepaper.gopluslabs.io/goplus-network/goplus-open-research/phishing-address-detection}. Accessed: April 20, 2026.

\bibitem{beery2025simulation}
T. Be'ery, ``Transaction Simulation White Paper,'' Zengo, Apr. 6, 2025. Available: \url{https://zengo.com/transaction-simulation-white-paper-ethereum-foundation-grant/}. Accessed: April 20, 2026.

\bibitem{kanezashi2022hetero}
H. Kanezashi, T. Suzumura, X. Liu, and T. Hirofuchi, ``Ethereum Fraud Detection with Heterogeneous Graph Neural Networks,'' arXiv:2203.12363, 2022. Available: \url{https://arxiv.org/abs/2203.12363}. Accessed: April 20, 2026.

\end{thebibliography}

\end{document}